% GENERATED FILE. Edit canonical lessons, then run npm run generate:books.
% canonical-source-hash: fnv1a64:e8141c851575c31a
% canonical-lessons: ZH-C488-gongchengshi, ZH-C488-chushi, ZH-C488-kexuejia, ZH-C488-fanyi, ZH-C488-zhuan

\chapter{What Do You Want to Be}
\label{ch:zh-what-do-you-want-to-be}
\hlchaptermodality{488}

\begin{chapteropening}
\textbf{By the end of this chapter:} \emph{I can say what I want to be: an engineer, a chef or a scientist, and talk about translating and earning.}

Complete the last lesson of chapter 488: I can say what I want to be: an engineer, a chef or a scientist, and talk about translating and earning.
\end{chapteropening}

\section[gōngchéngshī]{gōngchéngshī — engineer}

\label{lesson:ZH-C488-gongchengshi}



\begin{quote}
Before the new one: say the Mandarin for to be out of work, then the Mandarin for a factory.
\end{quote}

\subsection*{You'll want to know: gōngchéngshī}
\emph{gōngchéngshī} — \textquotedblleft{}engineer\textquotedblright{}.

\begin{quote}
\emph{wǒ yǐhòu xiǎng dāng gōngchéngshī} — I want to be an engineer.
\end{quote}

Its characters come later; for now it is a word for the ear.

\begin{cousinweb}
\emph{gōngchéngshī} is built from \zh{工} 'work', \emph{chéng} 'a measure' and \emph{shī} 'a master', all inherited from Middle Chinese, \emph{kuwng}, \emph{drjeng} and \emph{srij} (Guangyun, in Baxter's transcription), which glosses the second 'a rule' and the third 'a model to follow'; the sense 'engineering' is modern and not given there.
\end{cousinweb}

\begin{grammarlens}[title={What you've built}]
One more word for saying how it was.
\end{grammarlens}

\subsection*{Your turn}
\begin{itemize}
\raggedright
  \item \emph{Say it:} \emph{gōngchéngshī}
  \item \emph{Say it:} \emph{gōngchéngshī}, once more
  \item \emph{Say it:} \emph{wǒ yǐhòu xiǎng dāng gōngchéngshī} — I want to be an engineer.
  \item \emph{Recall:} say the Mandarin for to be out of work, then the Mandarin for a factory, then say \emph{gōngchéngshī} again
\end{itemize}

\begin{tcolorbox}[breakable,colback=teal!4,colframe=teal!35!black,title={Before you move on}]
What does gōngchéngshī mean? (\textquotedblleft{}Engineer\textquotedblright{}.) Say it once more.
\end{tcolorbox}
\section[chúshī]{chúshī — chef; professional cook}

\label{lesson:ZH-C488-chushi}



\begin{quote}
Before the new one: say the Mandarin for a factory, then the Mandarin for an engineer.
\end{quote}

\subsection*{You'll want to know: chúshī}
\emph{chúshī} — \textquotedblleft{}chef; professional cook\textquotedblright{}.

\begin{quote}
\emph{tā dāng le shí nián chúshī měi tiān dōu hěn píláo} — He was a chef for ten years and was worn out every day.
\end{quote}

Its characters come later; for now it is a word for the ear.

\begin{cousinweb}
\emph{chúshī} is built from \emph{chú} 'a kitchen' and \emph{shī} 'a master', both inherited from Middle Chinese, \emph{drju} and \emph{srij} (Guangyun, in Baxter's transcription), which glosses the first 'a kitchen' and the second 'a model to follow'.
\end{cousinweb}

\begin{grammarlens}[title={What you've built}]
One more word for saying how it was.
\end{grammarlens}

\subsection*{Your turn}
\begin{itemize}
\raggedright
  \item \emph{Say it:} \emph{chúshī}
  \item \emph{Say it:} \emph{chúshī}, once more
  \item \emph{Say it:} \emph{tā dāng le shí nián chúshī měi tiān dōu hěn píláo} — He was a chef for ten years and was worn out every day.
  \item \emph{Recall:} say the Mandarin for a factory, then the Mandarin for an engineer, then say \emph{chúshī} again
\end{itemize}

\begin{tcolorbox}[breakable,colback=teal!4,colframe=teal!35!black,title={Before you move on}]
What does chúshī mean? (\textquotedblleft{}Chef; professional cook\textquotedblright{}.) Say it once more.
\end{tcolorbox}
\section[kēxuéjiā]{kēxuéjiā — scientist}

\label{lesson:ZH-C488-kexuejia}



\begin{quote}
Before the new one: say the Mandarin for an engineer, then the Mandarin for a chef.
\end{quote}

\subsection*{You'll want to know: kēxuéjiā}
\emph{kēxuéjiā} — \textquotedblleft{}scientist\textquotedblright{}.

\begin{quote}
\emph{xiǎoshíhou wǒ zuì xiǎng dāng kēxuéjiā} — When I was little, what I wanted most was to be a scientist.
\end{quote}

Its characters come later; for now it is a word for the ear.

\begin{cousinweb}
\emph{kēxuéjiā} is built from \emph{kē} 'a grade', \zh{学} 'to learn' and \zh{家} 'a household', all inherited from Middle Chinese, \emph{khwa}, \emph{haewk} and \emph{kae} (Guangyun, in Baxter's transcription), which glosses the first 'a rule' and 'a class'. The senses 'science' and 'someone skilled in it' are modern and not given there.
\end{cousinweb}

\begin{grammarlens}[title={What you've built}]
One more word for saying how it was.
\end{grammarlens}

\subsection*{Your turn}
\begin{itemize}
\raggedright
  \item \emph{Say it:} \emph{kēxuéjiā}
  \item \emph{Say it:} \emph{kēxuéjiā}, once more
  \item \emph{Say it:} \emph{xiǎoshíhou wǒ zuì xiǎng dāng kēxuéjiā} — When I was little, what I wanted most was to be a scientist.
  \item \emph{Recall:} say the Mandarin for an engineer, then the Mandarin for a chef, then say \emph{kēxuéjiā} again
\end{itemize}

\begin{tcolorbox}[breakable,colback=teal!4,colframe=teal!35!black,title={Before you move on}]
What does kēxuéjiā mean? (\textquotedblleft{}Scientist\textquotedblright{}.) Say it once more.
\end{tcolorbox}
\section[fānyì]{fānyì — to translate; to interpret}

\label{lesson:ZH-C488-fanyi}



\begin{quote}
Before the new one: say the Mandarin for a chef, then the Mandarin for a scientist.
\end{quote}

\subsection*{You'll want to know: fānyì}
\emph{fānyì} — \textquotedblleft{}to translate; to interpret\textquotedblright{}.

\begin{quote}
\emph{nǐ néng bāng wǒ fānyì zhège ma} — Can you translate this for me?
\end{quote}

Its characters come later; for now it is a word for the ear.

\begin{cousinweb}
\emph{fānyì} is built from \emph{fān} 'to turn over' and \emph{yì} 'to translate', both inherited from Middle Chinese, \emph{phjon} and \emph{yek} (Guangyun, in Baxter's transcription), which glosses the first 'to overturn' and the second 'to pass words on', the name of the interpreters for the peoples of the north.
\end{cousinweb}

\begin{grammarlens}[title={What you've built}]
One more word for saying how it was.
\end{grammarlens}

\subsection*{Your turn}
\begin{itemize}
\raggedright
  \item \emph{Say it:} \emph{fānyì}
  \item \emph{Say it:} \emph{fānyì}, once more
  \item \emph{Say it:} \emph{nǐ néng bāng wǒ fānyì zhège ma} — Can you translate this for me?
  \item \emph{Recall:} say the Mandarin for a chef, then the Mandarin for a scientist, then say \emph{fānyì} again
\end{itemize}

\begin{tcolorbox}[breakable,colback=teal!4,colframe=teal!35!black,title={Before you move on}]
What does fānyì mean? (\textquotedblleft{}To translate; to interpret\textquotedblright{}.) Say it once more.
\end{tcolorbox}
\section[zhuàn]{zhuàn — to earn money; to make a profit}

\label{lesson:ZH-C488-zhuan}



\begin{quote}
Before the new one: say the Mandarin for a scientist, then the Mandarin for to translate.
\end{quote}

\subsection*{You'll want to know: zhuàn}
\emph{zhuàn} — \textquotedblleft{}to earn money; to make a profit\textquotedblright{}.

\begin{quote}
\emph{tā yí ge yuè zhuàn yíwàn kuài} — He earns ten thousand yuan a month.
\end{quote}

Its characters come later; for now it is a word for the ear.

\begin{cousinweb}
\emph{zhuàn} has no entry in the Guangyun, the dictionary this track cites for older forms. No dictionary we use records an older form, so its history is not settled.
\end{cousinweb}

\begin{grammarlens}[title={What you've built}]
One more word for saying how it was.
\end{grammarlens}

\subsection*{Your turn}
\begin{itemize}
\raggedright
  \item \emph{Say it:} \emph{zhuàn}
  \item \emph{Say it:} \emph{zhuàn}, once more
  \item \emph{Say it:} \emph{tā yí ge yuè zhuàn yíwàn kuài} — He earns ten thousand yuan a month.
  \item \emph{Recall:} say the Mandarin for a scientist, then the Mandarin for to translate, then say \emph{zhuàn} again
\end{itemize}

\begin{tcolorbox}[breakable,colback=teal!4,colframe=teal!35!black,title={Before you move on}]
What does zhuàn mean? (\textquotedblleft{}To earn money; to make a profit\textquotedblright{}.) Say it once more.
\end{tcolorbox}
